\section{Explicit finite premises}\label{app:finite-data}

\subsection{Matrix codes and representative planes}
For $0\leq v<512$ define $\operatorname{mat}(v)_{ij}$ to be the bit
at position $3i+j$ of $v$. Thus the least significant bit represents
$E_{00}$, and code $272$ represents $\operatorname{diag}(0,1,1)$.
A displayed tuple of codes denotes their $\F$-linear span. The frozen
RREF convention chooses pivots from bit $8$ down to bit $0$.
The following ranks are the sorted ranks of the three nonzero elements,
not necessarily the ranks in basis order.

\begin{table}[htbp]
\centering
\caption{The eight all-high planes. Every inherited label is $18$;
every improved quotient lower bound is $19$.}
\label{tab:planes}
\begin{tabular}{rcrr}
\toprule
Index & Basis codes & Rank multiset & Orbit size\\
\midrule
484 & $(19,10)$ & $(2,2,2)$ & 98\\
485 & $(20,10)$ & $(2,2,2)$ & $2{,}352$\\
486 & $(68,10)$ & $(2,2,2)$ & $1{,}176$\\
487 & $(84,10)$ & $(2,3,3)$ & $3{,}528$\\
488 & $(96,10)$ & $(2,2,3)$ & $14{,}112$\\
489 & $(258,10)$ & $(2,2,2)$ & $4{,}704$\\
490 & $(275,10)$ & $(2,3,3)$ & $4{,}704$\\
491 & $(163,84)$ & $(3,3,3)$ & $1{,}344$\\
\bottomrule
\end{tabular}
\end{table}

Their orbit sizes sum to $32{,}018$. There are $43{,}435$ two-planes
in total. The six other orbit representatives have bases
$(2,1)$, $(10,1)$, $(16,1)$, $(20,1)$, $(84,1)$, and $(160,1)$,
with indices $478$--$483$. Each contains a rank-one matrix and is
outside the hypothesis of Corollary~\ref{cor:raises}.

The total number of subspaces is obtained by counting ordered bases:
each $d$-space has $\prod_{j=0}^{d-1}(2^d-2^j)$ ordered bases,
whereas $V$ has $\prod_{j=0}^{d-1}(2^9-2^j)$ independent
ordered $d$-tuples. Thus
\[
 \sum_{d=0}^9 \qbinom{9}{d}=8{,}283{,}458,
 \qquad
 \qbinom{n}{d}=\prod_{j=0}^{d-1}\frac{2^n-2^j}{2^d-2^j}.
\]
Similarly $\qbinom{9}{2}=43{,}435$, and
$\sum_{d=1}^6\qbinom{7}{d}=29{,}210$. These count actual subspaces;
the formal orbit classification supplies the separate coverage claim.

\subsection{The normalized hyperplanes}
The spaces $S=(272,4,2,1)$ and $R=(4,2,1)$ have dimensions four
and three. Their frozen values are $14$ and $15$. The other fourteen
hyperplanes of $S$ are listed below; each has frozen value $17$.
The index is the suffix in the formal name \code{affineHyperplane}.
\begin{center}
\begin{tabular}{rc@{\qquad}rc}
\toprule
Index & Basis codes & Index & Basis codes\\
\midrule
0 & $(272,4,2)$ & 7 & $(276,2,1)$\\
1 & $(273,4,2)$ & 8 & $(272,5,2)$\\
2 & $(272,4,1)$ & 9 & $(273,5,2)$\\
3 & $(274,4,1)$ & 10 & $(272,6,1)$\\
4 & $(272,4,3)$ & 11 & $(274,6,1)$\\
5 & $(273,4,3)$ & 12 & $(272,5,3)$\\
6 & $(272,2,1)$ & 13 & $(273,5,3)$\\
\bottomrule
\end{tabular}
\end{center}
Equivalently these are
$\{\alpha p+e_0w^T:a\cdot w=b\alpha\}$ for
$a\in\F^3\setminus0$ and $b\in\F$. This formula proves that
the list exhausts the fourteen affine-plane spans used in the proof.

The line representatives in the small premise interface are
$(1)$, $(17)$, and $(273)$, of matrix ranks one, two, and three.
They are respectively equivalent to the frozen line representatives
$(1)$, $(10)$, and $(84)$, all with label $19$.

\subsection{Integer branch-certificate sizes}
Table~\ref{tab:branches} describes the Lean integer branch certificates.
A zero-forcing row has target-level
bound $18$ and removes its quotient directions by nonnegativity.
Source rows retain their proved labels in~\eqref{eq:retained}.
All systems begin with $127$ nonzero directions and $29{,}210$ rows;
retaining fewer rows is a weakening, so infeasibility of the retained
system proves infeasibility of the full one. The coefficient tests,
row bindings, and exhaustive branches are part of the formal proof.
% Generated statistics of the integer branch certificates, not DRAT.
\begin{table}[htbp]
\centering
\caption{Integer branch certificates after zero-weight elimination.}
\label{tab:branches}
\begin{tabular}{rrrrr}
\toprule
Plane & Live variables & Source rows & Zero-forcing rows & Leaves\\
\midrule
484 & 43 & 1{,}665 & 84 & 569\\
485 & 41 & 421 & 86 & 40\\
486 & 41 & 262 & 86 & 14\\
487 & 48 & 668 & 79 & 89\\
488 & 41 & 413 & 86 & 32\\
489 & 32 & 371 & 95 & 43\\
490 & 43 & 713 & 84 & 115\\
491 & 49 & 726 & 78 & 120\\
\bottomrule
\end{tabular}
\end{table}


\subsection{Substitution and a source bound of twelve}
\label{app:source-substitution}
The substitution rule used for the source bounds has a useful
formulation in terms of slices. Let $D$ be a finite-dimensional
$\F$-vector space, let $\beta,\gamma$ be finite index sets, and write
a tensor as $S:\beta\times\gamma\to D$. Its second-factor slices
are the functions $S_b:c\mapsto S(b,c)$.

\begin{lemma}[Substitution with a surviving flattening]
\label{lem:source-substitution}
Choose $b_1,\ldots,b_k\in\beta$ such that the slices $S_{b_l}$
are linearly independent, and $c_1,\ldots,c_f\in\gamma$.
Suppose that for every nonzero $y\in\F^f$ there exist $b\in\beta$
and $\phi\in D^*$ satisfying
\begin{equation}\label{eq:substitution-witness}
 \phi\left(\sum_j y_j S(b_l,c_j)\right)=0\quad(1\leq l\leq k),
 \qquad
 \phi\left(\sum_j y_j S(b,c_j)\right)\ne0.
\end{equation}
Then $R_{\F}(S)\geq k+f$.
\end{lemma}
\begin{proof}
Consider a decomposition $S(b,c)=\sum_{t=1}^r a_t B_t(b)C_t(c)$.
If $k>0$, the nonzero slice at $b_1$ forces $B_{t_0}(b_1)=1$
for some $t_0$.
Replace the slice family by
\[
 S'(b,c)=S(b,c)+B_{t_0}(b)S(b_1,c).
\]
The $t_0$-th term cancels, leaving a decomposition with at most $r-1$
terms. The other selected slices remain independent: a relation among
them would give a relation among the original $k$ selected slices,
with the coefficient of $S_{b_1}$ adjusted accordingly.
For a witness in~\eqref{eq:substitution-witness}, the added multiple
of $S_{b_1}$ vanishes under $\phi$. Hence the same condition holds
for $S'$ and its $k-1$ remaining selected slices.

Iterate $k$ times. In the final family, no nonzero linear combination
of the $f$ selected columns vanishes, since the surviving witness
detects it at some $b$. These columns are therefore independent in
the flattening with column index $c$ and row space $D\otimes\F^\beta$.
Its rank is at least $f$, whereas a decomposition with at most $r-k$
terms gives rank at most $r-k$. Thus $r\geq k+f$.
\end{proof}

For a source actually used in the formal development, set
\[
 J=\begin{pmatrix}0&0&1\\0&1&0\\1&0&0\end{pmatrix},
 \qquad K=E_{22},\qquad W=H_J\cap H_K.
\]
This is $W_{11}$ in the catalogue. Quotient coordinates are
$X\mapsto(\langle J,X\rangle,\langle K,X\rangle)$, so its slice
family, with $b=(j,k)$ and $c=(i,l)$, is
\begin{equation}\label{eq:source-twelve-slices}
 S((j,k),(i,l))=
    (\delta_{i+j,2}\delta_{kl},\ \delta_{i,2}\delta_{j,2}\delta_{kl}).
\end{equation}
Choose the six slices with $j=0,1$ and all $k$, and the six columns
with $i=0,2$ and all $l$. The slices are independent: in the first
component, each has a different nonzero coordinate $(i,l)=(2-j,k)$.

Write a combination of the selected columns as $y=(a,b)\in\F^3\oplus\F^3$,
where $a$ corresponds to $i=0$ and $b$ to $i=2$.
If $b\ne0$, the second coordinate functional annihilates all six
selected slice evaluations and detects $b_k\ne0$ in a slice $(2,k)$.
If $b=0$ and $a\ne0$, the first coordinate functional instead
annihilates those evaluations and detects $a_k\ne0$ in a slice
$(2,k)$. This proves~\eqref{eq:substitution-witness} for every
nonzero $y$. Lemma~\ref{lem:source-substitution} yields
$R_{\F}(T_{W_{11}})\geq6+6=12$.
This example and~\eqref{eq:three-hyperplane-bound} illustrate the two
complementary operations used in the source proofs: preserving a
flattening through substitution, and combining quotient capacities.
